% ROZDZIAŁ 8 — do włączenia w main.tex
\chapter{Twierdzenia graniczne dla całki Lebesgue'a}
\label{ch:twierdzenia-graniczne}

\section{Dlaczego nie wolno bezwarunkowo przestawiać granicy i całki?}

Punktem wyjścia jest pytanie, kiedy dla $f_n\to f$ wolno
zamienić kolejność granicy i całkowania:
\[
\int_X f(x)\,\mathrm d\mu(x)
=\lim_{n\to\infty}\int_X f_n(x)\,\mathrm d\mu(x).
\]
Zbieżność punktowa sama nie wystarcza. Na przykład na
$(0,1)$ z miarą Lebesgue'a funkcje
$f_n=n\chi_{(0,1/n)}$ zbiegają do zera w każdym
punkcie, ale ich całki wynoszą $1$.
Twierdzenia tego rozdziału podają precyzyjne warunki,
w których przejście do granicy jest uprawnione.

W całym rozdziale $(X,\mathcal F,\mu)$ jest przestrzenią
z miarą. Każde występujące całkowanie oznacza całkowanie
w sensie rozdziału \ref{ch:calkowanie-lebesgue}.
Nie zakładamy, że $\mu$ jest skończona ani zupełna,
jeśli nie zaznaczamy tego wyraźnie. Symbol $f_n\uparrow f$
oznacza niemalejącą zbieżność punktową do $f$.

\section{Fundamentalny lemat aproksymacyjny}

\begin{lemat}[Zbieżność całek prostej funkcji na zbiorach rosnących]
\label{lem:simple-exhaustion}
Niech $s:X\to[0,\infty)$ będzie mierzalną funkcją
prostą, a $(E_n)_{n\ge1}\subseteq\mathcal F$
niemalejącym ciągiem zbiorów. Jeśli
$E=\bigcup_{n\ge1}E_n$, to
\[
\lim_{n\to\infty}\int_{E_n}s(x)\,\mathrm d\mu(x)
=\int_Es(x)\,\mathrm d\mu(x).
\]
Równość zachodzi również wtedy, gdy prawa strona
wynosi $+\infty$.
\end{lemat}
\begin{proof}
Niech $s=\sum_{j=1}^m a_j\chi_{A_j}$
będzie reprezentacją na mierzalnym podziale, gdzie
$a_j\ge0$. Z definicji całki po zbiorze
$\int_{E_n}s\,\mathrm d\mu(x)=
\sum_{j=1}^m a_j\mu(A_j\cap E_n)$.
Zbiory $A_j\cap E_n$ rosną do $A_j\cap E$, więc
ciągłość miary z dołu (twierdzenie
\ref{thm:measure-continuity}) daje
$\mu(A_j\cap E_n)\uparrow\mu(A_j\cap E)$.
Ponieważ suma zawiera skończenie wiele nieujemnych
składników, możemy przejść do granicy w każdym z nich;
przy współczynniku $a_j=0$ obie strony są zerem.
Otrzymujemy wzór dla $E$.
\end{proof}

\section{Twierdzenie Beppo Leviego}

\begin{twierdzenie}[O zbieżności monotonicznej; Beppo Levi]
\label{thm:mct}
Niech $(f_n)_{n\ge1}$ będzie ciągiem mierzalnych funkcji
$f_n:X\to[0,\infty]$, takim że dla każdego $x\in X$
$f_n(x)\le f_{n+1}(x)$. Niech
$f(x)=\lim_{n\to\infty}f_n(x)=\sup_nf_n(x)$.
Wtedy $f$ jest mierzalna i
\begin{equation}\label{eq:mct}
\int_Xf(x)\,\mathrm d\mu(x)
=\lim_{n\to\infty}\int_Xf_n(x)\,\mathrm d\mu(x)
=\sup_{n\ge1}\int_Xf_n(x)\,\mathrm d\mu(x).
\end{equation}
Dopuszczamy wartość $+\infty$ i nie wymagamy
żadnego założenia o $\mu(X)$.
\end{twierdzenie}
\begin{proof}
Mierzalność $f$ wykazano w twierdzeniu
\ref{thm:pointwise-limits}. Z monotoniczności całki
(rozdział 7) mamy
$L:=\sup_n\int_Xf_n\,\mathrm d\mu(x)\le
\int_Xf\,\mathrm d\mu(x)$.
Trzeba wykazać trudniejszą nierówność odwrotną.

Ustalmy dowolną nieujemną mierzalną funkcję prostą
$s\le f$ oraz liczbę $\theta\in(0,1)$.
Definiujemy
$E_n=\{x\in X:f_n(x)\ge\theta s(x)\}$.
Zbiory $E_n$ są mierzalne i rosną.
Ponadto $\bigcup_nE_n=X$: jeżeli $s(x)=0$,
to $x\in E_n$ dla każdego $n$; jeżeli $s(x)>0$,
to $f(x)\ge s(x)>\theta s(x)$, więc z
$f_n(x)\uparrow f(x)$ wynika $f_n(x)\ge\theta s(x)$
dla dostatecznie dużego $n$. Argument działa także
gdy $f(x)=+\infty$, ponieważ $s(x)$ jest skończona.

Z nierówności $f_n\ge\theta s\chi_{E_n}$,
monotoniczności i jednorodności całki wynika
\[
\int_Xf_n\,\mathrm d\mu(x)\ge
\theta\int_{E_n}s(x)\,\mathrm d\mu(x).
\]
Przechodząc do granicy, stosując lemat
\ref{lem:simple-exhaustion} i $E_n\uparrow X$,
dostajemy $L\ge\theta\int_Xs\,\mathrm d\mu(x)$.
Niech $\theta\uparrow1$; dla całki skończonej
wynika $L\ge\int_Xs\,\mathrm d\mu(x)$, a jeśli
$\int_Xs=\infty$, to z powyższej nierówności
już dla jednego $\theta>0$ otrzymujemy $L=\infty$.
Na koniec bierzemy supremum po wszystkich
nieujemnych prostych $s\le f$. Na mocy definicji
\ref{def:integral-positive} daje to
$L\ge\int_Xf\,\mathrm d\mu(x)$. Obie nierówności
kończą dowód.
\end{proof}

\begin{wniosek}[Konstruktywne obliczanie całki]
\label{cor:simple-approx-integral}
Jeżeli $f:X\to[0,\infty]$ jest mierzalna, a
$(s_n)$ jest \emph{dowolnym} ciągiem nieujemnych
funkcji prostych mierzalnych takim, że $s_n\uparrow f$
wszędzie, to
$\int_Xf\,\mathrm d\mu(x)=
\lim_{n\to\infty}\int_Xs_n\,\mathrm d\mu(x)$.
W szczególności można użyć funkcji z
\eqref{eq:canonical-simple}.
\end{wniosek}
\begin{proof}
To bezpośrednie zastosowanie twierdzenia
\ref{thm:mct} do ciągu $s_n$; całki funkcji prostych
znamy już z definicji \ref{def:integral-simple}.
\end{proof}

\begin{przyklad}[Ciąg monotoniczny o nieskończonej granicy całek]
Na $(0,1)$ z miarą Lebesgue'a funkcje
$f_n(x)=\min\{n,1/x\}$ są niemalejące i zbiegają
do $1/x$. Twierdzenie \ref{thm:mct} daje
$\int_{(0,1)}x^{-1}\,\mathrm d\lambda_1(x)=
\lim_n\int_{(0,1)}f_n\,\mathrm d\lambda_1(x)=+\infty$,
co jest zgodne z bezpośrednim rachunkiem z
rozdziału \ref{ch:calkowanie-lebesgue}.
Nie wolno wywnioskować, że same $f_n$ mają
nieograniczone wartości punktowe: każda z nich
jest ograniczona przez $n$.
\end{przyklad}

\begin{twierdzenie}[Zbieżność monotoniczna prawie wszędzie]
\label{thm:mct-ae}
Niech $f_n,f:X\to[0,\infty]$ będą mierzalne.
Jeśli $f_n\le f_{n+1}$ i $f_n\to f$
$\mu$-prawie wszędzie, to zachodzi równość
\eqref{eq:mct}.
\end{twierdzenie}
\begin{proof}
Niech $N\in\mathcal F$ będzie zbiorem miary zero,
poza którym monotoniczność i zbieżność zachodzą.
Funkcje $g_n=f_n\chi_{X\setminus N}$
są mierzalne, nieujemne i niemalejące wszędzie,
a ich granica jest równa
$g=f\chi_{X\setminus N}$.
Stosujemy twierdzenie \ref{thm:mct} do $g_n$.
Równość całek $g_n$ i $f_n$ oraz całek $g$ i $f$
wynika z twierdzenia \ref{thm:ae-integral-invariance}.
Dzięki temu zastosowanie wersji punktowej
Beppo Leviego daje dokładnie żądaną równość.
\end{proof}

\section{Addytywność i liniowość całki}

\begin{twierdzenie}[Addytywność całki nieujemnej]
\label{thm:positive-additivity}
Dla dowolnych mierzalnych $f,g:X\to[0,\infty]$
zachodzi
\begin{equation}\label{eq:positive-additivity}
\int_X(f+g)(x)\,\mathrm d\mu(x)
=\int_Xf(x)\,\mathrm d\mu(x)+
\int_Xg(x)\,\mathrm d\mu(x).
\end{equation}
Obie strony należą do $[0,\infty]$; dopuszczamy
$+\infty$ po którejkolwiek stronie.
\end{twierdzenie}
\begin{proof}
Z twierdzenia \ref{thm:simple-up} wybierzmy
nieujemne funkcje proste $s_n\uparrow f$ i
$t_n\uparrow g$. Wtedy $s_n+t_n\uparrow f+g$.
Twierdzenie \ref{thm:mct}, a następnie addytywność
całki prostej, dają
\[
\int_X(f+g)\,\mathrm d\mu(x)
=\lim_n\int_X(s_n+t_n)\,\mathrm d\mu(x)
=\lim_n\left(\int_Xs_n\,\mathrm d\mu(x)+
\int_Xt_n\,\mathrm d\mu(x)\right).
\]
Ponowne zastosowanie twierdzenia \ref{thm:mct}
do $s_n$ i $t_n$ daje tezę; suma dwóch ciągów
niemalejących liczb nieujemnych ma granicę równą
sumie ich granic także w $[0,\infty]$.
\end{proof}

\begin{wniosek}[Rozkład całki na rozłączne zbiory]
\label{cor:integral-countable-partition}
Niech $f:X\to[0,\infty]$ będzie mierzalna, a
$(A_n)_{n\ge1}$ przeliczalną rodziną parami rozłącznych
zbiorów z $\mathcal F$. Wtedy
\begin{equation}\label{eq:integral-partition}
\int_{\bigcup_{n\ge1}A_n}f(x)\,\mathrm d\mu(x)
=\sum_{n=1}^{\infty}\int_{A_n}f(x)\,\mathrm d\mu(x).
\end{equation}
\end{wniosek}
\begin{proof}
Dla każdego $N$ zachodzi punktowo
$f\chi_{\cup_{n=1}^NA_n}=
\sum_{n=1}^Nf\chi_{A_n}$.
Addytywność całki daje wzór dla sumy skończonej.
Gdy $N\to\infty$, lewe funkcje rosną do
$f\chi_{\cup_{n\ge1}A_n}$. Zastosowanie
Beppo Leviego kończy dowód.
\end{proof}

\begin{twierdzenie}[Liniowość całki funkcji całkowalnych]
\label{thm:integral-linear-real}
Jeżeli $f,g:X\to\mathbb R$ są mierzalne i
$\int_X|f|\,\mathrm d\mu(x)<\infty$,
$\int_X|g|\,\mathrm d\mu(x)<\infty$, to dla
$a,b\in\mathbb R$ funkcja $af+bg$ jest całkowalna i
\[
\int_X(af+bg)(x)\,\mathrm d\mu(x)
=a\int_Xf(x)\,\mathrm d\mu(x)+
b\int_Xg(x)\,\mathrm d\mu(x).
\]
Ponadto $\bigl|\int_Xf\,\mathrm d\mu(x)\bigr|
\le\int_X|f|\,\mathrm d\mu(x)$.
\end{twierdzenie}
\begin{proof}
Dla $a\ge0$ rozkłady $(af)^+=af^+$ i
$(af)^-=af^-$ oraz jednorodność całki nieujemnej
dają $\int_X af=a\int_X f$. Dla $a<0$ korzystamy z
$(-f)^+=f^-$ i $(-f)^-=f^+$.

Z nierówności $|f+g|\le|f|+|g|$ i addytywności
nieujemnej wynika całkowalność sumy.
W każdym punkcie zachodzi tożsamość
\[
(f+g)^++f^-+g^-=(f+g)^-+f^++g^+.
\]
Wszystkie wyrazy mają skończone całki, bo są
ograniczone przez sumy $|f|+|g|$.
Całkujemy tę równość, stosujemy addytywność
nieujemną i przenosimy skończone składniki na
drugą stronę. Otrzymujemy
$\int_X(f+g)=\int_X f+\int_X g$.
Połączenie obu własności daje pełną liniowość.
Na koniec z $0\le f^\pm\le|f|$ i
$|\int_X f^+-\int_X f^-|\le\int_X f^++\int_X f^-
=\int_X|f|$ wynika oszacowanie.
\end{proof}

\begin{wniosek}[Nierówność trójkąta i struktura $L^1$]
\label{cor:integral-triangle}
Jeśli $f,g:X\to\mathbb R$ są całkowalne, to
$\int_X|f+g|\,\mathrm d\mu(x)\le
\int_X|f|\,\mathrm d\mu(x)+
\int_X|g|\,\mathrm d\mu(x)$.
Funkcjonał $f\mapsto\int_Xf\,\mathrm d\mu(x)$
jest liniowy na przestrzeni rzeczywistych
funkcji całkowalnych, a nierówność
$|\int_X f|\le\int_X|f|$ wyznacza jego ograniczenie.
Po utożsamieniu funkcji równych prawie wszędzie
otrzymamy przestrzeń normowaną $L^1$ w rozdziale 14.
\end{wniosek}
\begin{proof}
Całkujemy punktową nierówność $|f+g|\le|f|+|g|$
i stosujemy addytywność całki nieujemnej. Liniowość
i oszacowanie funkcjonału zostały udowodnione
w twierdzeniu \ref{thm:integral-linear-real}.
\end{proof}

\begin{definicja}[Całka funkcji zespolonej]
\label{def:complex-integral}
Funkcję mierzalną $f:X\to\mathbb C$ nazywamy
całkowalną, jeżeli $\int_X|f|\,\mathrm d\mu(x)<\infty$.
Wówczas określamy
$\int_Xf\,\mathrm d\mu(x)=
\int_X\operatorname{Re}f\,\mathrm d\mu(x)+
 i\int_X\operatorname{Im}f\,\mathrm d\mu(x)$.
\end{definicja}

\begin{twierdzenie}[Liniowość zespolona i oszacowanie modułu]
\label{thm:integral-complex-linear}
Jeśli $f,g:X\to\mathbb C$ są całkowalne i
$\alpha,\beta\in\mathbb C$, to $\alpha f+\beta g$
jest całkowalna,
$\int_X(\alpha f+\beta g)\,\mathrm d\mu(x)=
\alpha\int_Xf\,\mathrm d\mu(x)+
\beta\int_Xg\,\mathrm d\mu(x)$, a także
\[
\left|\int_Xf(x)\,\mathrm d\mu(x)\right|
\le\int_X|f(x)|\,\mathrm d\mu(x).
\]
\end{twierdzenie}
\begin{proof}
Z $|\alpha f+\beta g|\le|\alpha||f|+|\beta||g|$
wynika całkowalność. Liniowość otrzymujemy,
rozwijając części rzeczywiste i urojone i stosując
twierdzenie \ref{thm:integral-linear-real}.
Jeśli $I=\int_Xf\,\mathrm d\mu(x)\ne0$, wybierzmy
$\zeta=\overline I/|I|$, więc $|\zeta|=1$ i
$\zeta I=|I|$. Wtedy
$|I|=\int_X\operatorname{Re}(\zeta f)\,\mathrm d\mu(x)
\le\int_X|\operatorname{Re}(\zeta f)|\,\mathrm d\mu(x)
\le\int_X|f|\,\mathrm d\mu(x)$.
Dla $I=0$ nierówność jest oczywista.
\end{proof}

\section{Lemat Fatou}

\begin{twierdzenie}[Lemat Fatou]
\label{thm:fatou}
Niech $f_n:X\to[0,\infty]$ będą funkcjami mierzalnymi.
Wtedy
\begin{equation}\label{eq:fatou}
\int_X\liminf_{n\to\infty}f_n(x)\,\mathrm d\mu(x)
\le\liminf_{n\to\infty}\int_Xf_n(x)\,\mathrm d\mu(x).
\end{equation}
Nie wymagamy monotoniczności ciągu ani skończoności miary.
\end{twierdzenie}
\begin{proof}
Dla $m\ge1$ połóżmy
$h_m(x)=\inf_{n\ge m}f_n(x)$.
Na mocy twierdzenia \ref{thm:pointwise-limits}
każda $h_m$ jest mierzalna. Ponieważ z rosnącym $m$
infimum obejmuje coraz mniej wyrazów, funkcje $h_m$
są niemalejące. Z definicji granicy dolnej
$h_m(x)\uparrow\liminf_nf_n(x)$.
Twierdzenie \ref{thm:mct} daje
\[
\int_X\liminf_nf_n\,\mathrm d\mu(x)
=\lim_{m\to\infty}\int_Xh_m\,\mathrm d\mu(x).
\]
Dla wszystkich $n\ge m$ mamy $h_m\le f_n$, więc
$\int_Xh_m\,\mathrm d\mu(x)\le
\inf_{n\ge m}\int_Xf_n\,\mathrm d\mu(x)$.
Przechodzimy do granicy po $m$ i otrzymujemy
prawą stronę \eqref{eq:fatou}.
\end{proof}

\begin{przyklad}[Nierówność Fatou bywa ostra]
Na $\mathbb R$ z miarą Lebesgue'a niech
$f_n=\chi_{[n,n+1]}$. Dla każdego $x$ zachodzi
$f_n(x)\to0$, ale $\int_{\mathbb R} f_n(x)\,\mathrm d\lambda_1(x)=1$.
Zatem $\int_{\mathbb R}\liminf_nf_n(x)\,\mathrm d\lambda_1(x)=0$
i $\liminf_n\int_{\mathbb R} f_n(x)\,\mathrm d\lambda_1(x)=1$.
Masa może \emph{uciekać do nieskończoności}.
\end{przyklad}

\begin{wniosek}[Odwrotny lemat Fatou przy dominacji]
\label{cor:reverse-fatou}
Niech $f_n:X\to\mathbb R$ będą mierzalne i
$0\le f_n\le g$ $\mu$-prawie wszędzie, gdzie
$g:X\to[0,\infty]$ jest całkowalna. Wtedy
\[
\limsup_n\int_Xf_n\,\mathrm d\mu(x)
\le\int_X\limsup_nf_n\,\mathrm d\mu(x).
\]
\end{wniosek}
\begin{proof}
Ponieważ $g$ jest całkowalna, $g<\infty$ prawie
wszędzie. Usuwamy więc jeden wspólny mierzalny
zbiór zerowy, na którym może wystąpić $g=\infty$
lub naruszenie nierówności $0\le f_n\le g$.
Na tym zbiorze zastępujemy $f_n$ i $g$ zerem;
całki nie zmieniają się. Możemy teraz założyć,
że $g$ jest skończona wszędzie i $0\le f_n\le g$.
Funkcje $g-f_n$ są nieujemne i mierzalne, a
$\liminf_n(g-f_n)=g-\limsup_nf_n$ w punktach,
gdzie $g$ jest skończona. Ponieważ $g<\infty$
prawie wszędzie, lemat Fatou i liniowość całki
funkcji całkowalnych dają
$\int_X g-\int_X\limsup f_n\le
\int_X g-\limsup\int_X f_n$.
Odejmując skończoną liczbę $\int_X g$,
otrzymujemy tezę.
\end{proof}

\begin{twierdzenie}[Lemat Fatou z całkowalnym ograniczeniem od dołu]
\label{thm:fatou-integrable-lower-8}
Niech $f_n:X\to\overline{\mathbb R}$ będą mierzalne,
a $g:X\to[0,\infty]$ całkowalna. Załóżmy,
że $f_n\ge-g$ prawie wszędzie dla każdego $n$.
Wówczas każda całka $\int_Xf_n\,\mathrm d\mu(x)$
jest określona w $(-\infty,\infty]$, podobnie całka
funkcji $\liminf_nf_n$, oraz
\[
\int_X\liminf_{n\to\infty}f_n(x)\,\mathrm d\mu(x)
\le \liminf_{n\to\infty}\int_Xf_n(x)\,\mathrm d\mu(x).
\]
Nie jest wymagana skończoność miary $\mu(X)$.
\end{twierdzenie}
\begin{proof}
Bierzemy jeden wspólny zbiór wyjątkowy miary zero,
na którym mogą zawodzić nierówności $f_n\ge-g$
lub skończoność $g$. Zmieniamy tam wszystkie
funkcje na zero; ich całki pozostają niezmienione.
Możemy więc zakładać, że $g<\infty$ wszędzie oraz
$f_n+g\ge0$ dla każdego $n$. Lemat Fatou daje
\[
\int_X(\liminf_nf_n+g)\,\mathrm d\mu(x)
=\int_X\liminf_n(f_n+g)\,\mathrm d\mu(x)
\le\liminf_n\int_X(f_n+g)\,\mathrm d\mu(x).
\]
Aby uzasadnić rachunek całek bez odejmowania
nieskończoności, zauważmy, że $f_n^-\le g$,
a więc $\int_X f_n^-<\infty$. Punktowa tożsamość
$(f_n+g)+f_n^-=f_n^++g$ i addytywność całki
nieujemnej dają
$\int_X(f_n+g)+\int_Xf_n^-=
\int_Xf_n^++\int_Xg$.
Wolno odjąć skończoną całkę $f_n^-$ i otrzymać
$\int_X(f_n+g)=\int_Xf_n+\int_Xg$.
Analogicznie postępujemy z $\liminf_nf_n$.
Odejmujemy tę samą skończoną liczbę $\int_Xg$
i otrzymujemy tezę. Żadne wyrażenie typu
$\infty-\infty$ nie występuje.
\end{proof}

\begin{przyklad}[Brak całkowalnego ograniczenia od dołu]
Na $(0,1)$ połóżmy $f_n=-n\chi_{(0,1/n)}$.
Mamy $f_n(x)\to0$ dla każdego $x\in(0,1)$,
ale $\int_{(0,1)}f_n(x)\,\mathrm d\lambda_1(x)=-1$.
Zatem nierówność z twierdzenia
\ref{thm:fatou-integrable-lower-8} byłaby fałszywa,
gdyby usunąć warunek dominacji od dołu.
\end{przyklad}

\section{Twierdzenie Lebesgue'a o zbieżności zdominowanej}

\begin{twierdzenie}[O zbieżności zdominowanej; Lebesgue]
\label{thm:dct}
Niech $f_n:X\to\mathbb R$ będą mierzalne i
$f_n(x)\to f(x)$ dla $\mu$-prawie wszystkich $x$, gdzie
$f:X\to\mathbb R$ jest mierzalna. Załóżmy, że
istnieje nieujemna całkowalna funkcja mierzalna
$g:X\to[0,\infty]$ taka, że
$|f_n(x)|\le g(x)$ dla wszystkich $n$ i
$\mu$-prawie wszystkich $x$, przy czym zbiór
wyjątkowy można wybrać wspólny dla całego ciągu.
Wówczas $f_n$ i $f$ są całkowalne oraz
\begin{equation}\label{eq:dct}
\lim_{n\to\infty}\int_Xf_n(x)\,\mathrm d\mu(x)
=\int_Xf(x)\,\mathrm d\mu(x).
\end{equation}
\end{twierdzenie}
\begin{proof}
Ponieważ $g$ jest całkowalna, zbiór $\{g=\infty\}$
jest zerowy. Po dołączeniu go do wspólnego zbioru
wyjątkowego możemy zastąpić na nim wszystkie
$f_n$, funkcję $f$ i $g$ zerem. Całki pozostają
niezmienione, a teraz dla każdego $x$ zachodzi
$f_n(x)\to f(x)$, $|f_n(x)|\le g(x)<\infty$.
W szczególności operacje $g\pm f_n$ są wszędzie
określone i nieujemne.
Z ciągłości wartości bezwzględnej wynika $|f|\le g$.
Monotoniczność daje
$\int_X|f_n|\le\int_X g<\infty$ oraz
$\int_X|f|\le\int_X g<\infty$.

Funkcje $g+f_n$ i $g-f_n$ są nieujemne.
Stosując do nich lemat Fatou, otrzymujemy
\[
\int_X(g+f)\,\mathrm d\mu(x)
\le\liminf_n\int_X(g+f_n)\,\mathrm d\mu(x)
=\int_Xg\,\mathrm d\mu(x)+
\liminf_n\int_Xf_n\,\mathrm d\mu(x),
\]
\[
\int_X(g-f)\,\mathrm d\mu(x)
\le\liminf_n\int_X(g-f_n)\,\mathrm d\mu(x)
=\int_Xg\,\mathrm d\mu(x)-
\limsup_n\int_Xf_n\,\mathrm d\mu(x).
\]
W drugim przejściu wszystkie całki $f_n$ są
skończone i jednostajnie ograniczone przez $\int_X g$.
Po skróceniu skończonej liczby $\int_X g$ pierwsza
nierówność daje
$\int_X f\le\liminf_n\int_X f_n$, a druga
$\limsup_n\int_X f_n\le\int_X f$.
Ponieważ $\liminf\le\limsup$, obie granice
są równe $\int_X f$.
\end{proof}

\begin{wniosek}[Zbieżność zdominowana w $L^1$]
\label{cor:dct-l1}
Przy założeniach twierdzenia \ref{thm:dct} zachodzi
\[
\lim_{n\to\infty}\int_X|f_n(x)-f(x)|\,\mathrm d\mu(x)=0.
\]
Dla funkcji zespolonych wniosek i twierdzenie
\ref{thm:dct} także zachodzą, jeśli zastąpić
wartość bezwzględną modułem.
\end{wniosek}
\begin{proof}
Mamy $|f_n-f|\to0$ prawie wszędzie oraz
$|f_n-f|\le2g$ prawie wszędzie.
Stosujemy twierdzenie \ref{thm:dct} do nieujemnych
funkcji $|f_n-f|$. Dla funkcji zespolonych
stosujemy ten sam argument do modułu, a z
$\bigl|\int_X(f_n-f)\bigr|\le\int_X|f_n-f|$
wynika zbieżność całek.
\end{proof}

\begin{wniosek}[Twierdzenie o zbieżności ograniczonej]
\label{cor:bct}
Niech $\mu(X)<\infty$ oraz $f_n,f:X\to\mathbb R$
będą mierzalne. Jeśli $f_n\to f$ prawie wszędzie
i istnieje $M<\infty$ takie, że $|f_n|\le M$
prawie wszędzie dla wszystkich $n$, to
$\int_X f_n\to\int_X f$ i $\int_X|f_n-f|\to0$.
\end{wniosek}
\begin{proof}
Funkcja stała $g\equiv M$ jest całkowalna,
gdyż $\int_X g\,\mathrm d\mu(x)=M\mu(X)<\infty$.
Stosujemy twierdzenie \ref{thm:dct} i wniosek
\ref{cor:dct-l1}.
\end{proof}

\begin{przyklad}[Dlaczego potrzebna jest dominacja]
Na $(0,1)$ weźmy
$f_n=n\chi_{(0,1/n)}$.
Dla każdego ustalonego $x>0$ zachodzi $f_n(x)\to0$,
lecz $\int_{(0,1)}f_n\,\mathrm d\lambda_1(x)=1$.
Gdyby istniała całkowalna $g\ge f_n$ dla wszystkich
$n$, wówczas twierdzenie o zbieżności zdominowanej
dawałoby sprzeczność. Problem nie leży zatem
w zbieżności punktowej, lecz w braku wspólnej
całkowalnej funkcji dominującej.
\end{przyklad}

\begin{przyklad}[Dlaczego skończoność miary ma znaczenie]
Na $\mathbb R$ funkcje $f_n=\chi_{[n,n+1]}$
są ograniczone wspólną stałą $1$, a mimo to
$\int_{\mathbb R}f_n(x)\,\mathrm d\lambda_1(x)=1\not\to0$, chociaż $f_n\to0$ wszędzie.
Stała $1$ nie jest całkowalna na $\mathbb R$,
więc twierdzenie \ref{cor:bct} nie ma zastosowania.
\end{przyklad}

\section{Szeregi funkcji i całek}

\begin{twierdzenie}[Szeregi nieujemnych funkcji mierzalnych]
\label{thm:sum-nonnegative-integrals}
Niech $f_n:X\to[0,\infty]$ będą mierzalne. Wtedy
\begin{equation}\label{eq:sum-nonnegative-integrals}
\int_X\sum_{n=1}^{\infty}f_n(x)\,\mathrm d\mu(x)
=\sum_{n=1}^{\infty}\int_Xf_n(x)\,\mathrm d\mu(x),
\end{equation}
gdzie obie strony mogą być nieskończone.
\end{twierdzenie}
\begin{proof}
Sumy częściowe $s_N=\sum_{n=1}^N f_n$ są
mierzalne, nieujemne i rosną do
$s=\sum_{n\ge1}f_n$. Z twierdzenia \ref{thm:mct}
$\int_Xs\,\mathrm d\mu(x)=\lim_N\int_Xs_N\,\mathrm d\mu(x)$.
Dla każdego $N$ stosujemy skończoną addytywność
całki nieujemnej z twierdzenia
\ref{thm:positive-additivity}; granica jest
z definicji sumą szeregu nieujemnego.
\end{proof}

\begin{twierdzenie}[Całkowanie szeregów absolutnie sumowalnych]
\label{thm:sum-absolutely-integrable}
Niech $f_n:X\to\mathbb C$ będą mierzalne i
całkowalne oraz niech
$\sum_{n\ge1}\int_X|f_n(x)|\,\mathrm d\mu(x)<\infty$.
Wówczas $\sum_nf_n(x)$ jest bezwzględnie zbieżny
$\mu$-prawie wszędzie. Jeśli $F$ oznacza jego sumę
na zbiorze zbieżności i zero na pozostałym zbiorze,
to $F$ jest całkowalna i
\begin{equation}\label{eq:interchange-series}
\int_XF(x)\,\mathrm d\mu(x)
=\sum_{n=1}^{\infty}\int_Xf_n(x)\,\mathrm d\mu(x).
\end{equation}
\end{twierdzenie}
\begin{proof}
Niech $G=\sum_{n\ge1}|f_n|$, mierzalna suma
nieujemna. Z twierdzenia
\ref{thm:sum-nonnegative-integrals}
$\int_XG\,\mathrm d\mu(x)=
\sum_n\int_X|f_n|\,\mathrm d\mu(x)<\infty$.
Na mocy twierdzenia \ref{thm:markov-nonnegative}
$G(x)<\infty$ prawie wszędzie. Poza zbiorem
zerowym szereg $\sum_nf_n(x)$ zbiega więc
bezwzględnie. Zbiór zbieżności jest mierzalny,
ponieważ zarówno części rzeczywiste, jak i
urojone sum częściowych są mierzalne, a warunek
Cauchy'ego można zapisać jako przeliczalną
kombinację mierzalnych nierówności.
Dla sum częściowych $S_N=\sum_{n=1}^Nf_n$
mamy $|S_N|\le G$ oraz $S_N\to F$ prawie wszędzie.
Zespolona wersja twierdzenia o zbieżności zdominowanej
daje $\int_XS_N\to\int_XF$. Liniowość całki
zespolonej pozwala zapisać
$\int_XS_N=\sum_{n=1}^N\int_Xf_n$, co daje tezę.
\end{proof}

\begin{twierdzenie}[Wzór warstwowy na całkę nieujemną]
\label{thm:layer-cake-8}
Niech $f:X\to[0,\infty]$ będzie mierzalna.
Wtedy funkcja $t\mapsto\mu(\{x:f(x)>t\})$
jest borelowska na $[0,\infty)$ i zachodzi
\begin{equation}\label{eq:layer-cake-8}
\int_Xf(x)\,\mathrm d\mu(x)
=\int_{[0,\infty)}\mu(\{x:f(x)>t\})\,
\mathrm d\lambda_1(t).
\end{equation}
Obie strony mogą być nieskończone. Dowód nie korzysta
z twierdzenia Tonellego, które pojawi się później.
\end{twierdzenie}
\begin{proof}
Funkcja $t\mapsto\mu(\{f>t\})$ jest \emph{nierosnąca},
a każda funkcja monotoniczna
o wartościach w $[0,\infty]$ jest borelowsko mierzalna.
Najpierw niech $s=\sum_{j=1}^m a_j\chi_{A_j}$
będzie nieujemną funkcją prostą, gdzie zbiory $A_j$
tworzą mierzalny podział $X$. Dla $t\ge0$
\[
\mu(\{s>t\})=\sum_{j=1}^m
\mu(A_j)\chi_{[0,a_j)}(t).
\]
Całkowanie skończonej sumy po $t$ daje
$\sum_ja_j\mu(A_j)=\int_Xs\,\mathrm d\mu(x)$;
przy $a_j=0$ stosujemy umowę $0\cdot\infty=0$.
Dla ogólnej $f$ weźmy $s_n\uparrow f$ z
twierdzenia~\ref{thm:simple-up}. Dla każdego $t\ge0$
mamy $\{s_n>t\}\uparrow\{f>t\}$. Ciągłość miary
z dołu daje
$\mu(\{s_n>t\})\uparrow\mu(\{f>t\})$.
Stosujemy Beppo Leviego na $[0,\infty)$ i na $X$:
\[
\int_{[0,\infty)}\mu(\{f>t\})\,
\mathrm d\lambda_1(t)
=\lim_n\int_Xs_n\,\mathrm d\mu(x)
=\int_Xf\,\mathrm d\mu(x).
\]
\end{proof}

\begin{wniosek}[Uogólniona formuła rozkładu na poziomice]\label{idx:08:wniosek:20}

Dla każdego $p>0$ i mierzalnej $f:X\to\mathbb R$
zachodzi
\[
\int_X|f(x)|^p\,\mathrm d\mu(x)
=p\int_{[0,\infty)}t^{p-1}
\mu(\{x:|f(x)|>t\})\,\mathrm d\lambda_1(t).
\]
\end{wniosek}
\begin{proof}
Niech najpierw $s=\sum_{j=1}^m a_j\chi_{A_j}$,
gdzie $a_j\ge0$, a $(A_j)$ jest mierzalnym podziałem.
Dla $t\ge0$ zachodzi
$\mu(\{s>t\})=\sum_j\mu(A_j)\chi_{[0,a_j)}(t)$.
Z addytywności całki nieujemnej i elementarnego wzoru
$\int_0^{a}pt^{p-1}\,\mathrm dt=a^p$ dla $a\ge0$
(w sensie całki Riemanna na przedziałach
$[\delta,a]$, a następnie przy $\delta\downarrow0$)
otrzymujemy
\[
p\int_{[0,\infty)}t^{p-1}\mu(\{s>t\})\,
\mathrm d\lambda_1(t)=\sum_j a_j^p\mu(A_j)
=\int_Xs(x)^p\,\mathrm d\mu(x).
\]
Przy $t=0$ i $p<1$ kładziemy wartość czynnika
$t^{p-1}$ równą zero; nie zmienia to całki.
Dla ogólnej funkcji $|f|$ wybieramy $s_n\uparrow|f|$
ze zbioru nieujemnych funkcji prostych.
Wtedy $s_n^p\uparrow|f|^p$ i
$\mu(\{s_n>t\})\uparrow\mu(\{|f|>t\})$.
Dwukrotne zastosowanie twierdzenia Beppo Leviego
daje żądaną równość bez używania twierdzenia Tonellego.
\end{proof}

\section{Przykłady obliczeń i kryteria całkowalności}

\begin{przyklad}[Całka potęgi blisko zera]
\label{ex:power-integral}
Niech $\alpha\in\mathbb R$ i
$f(x)=x^{-\alpha}$ na $(0,1]$.
Jeżeli $\alpha<1$, to funkcja jest całkowalna oraz
\begin{equation}\label{eq:power-integral}
\int_{(0,1]}x^{-\alpha}\,\mathrm d\lambda_1(x)
=\frac1{1-\alpha}.
\end{equation}
Jeśli $\alpha\ge1$, całka wynosi $+\infty$.
\end{przyklad}
\begin{proof}
Dla $\alpha\ge1$ mamy $x^{-\alpha}\ge x^{-1}$
na $(0,1]$, a rozbieżność drugiej całki
wykazaliśmy w rozdziale 7. Dla $\alpha<1$
rozważmy $f_n=f\chi_{[1/n,1]}$.
Wtedy $f_n\uparrow f$ na $(0,1]$.
Funkcja $f_n$ jest ograniczona i ciągła na
$[1/n,1]$, więc jej całka Lebesgue'a zgadza się
z całką Riemanna na tym przedziale (zob.
twierdzenie \ref{thm:riemann-to-lebesgue} poniżej).
Z rachunku elementarnego
$\int_{[1/n,1]}x^{-\alpha}\,\mathrm d\lambda_1(x)
=(1-n^{\alpha-1})/(1-\alpha)$.
Twierdzenie \ref{thm:mct} daje po przejściu do
granicy \eqref{eq:power-integral}.
\end{proof}

\begin{przyklad}[Kryterium logarytmiczne]
Dla $p\in\mathbb R$ rozważmy na $(0,e^{-1}]$
funkcję $f_p(x)=\bigl(x(\log(1/x))^p\bigr)^{-1}$.
Całka z $f_p$ jest skończona dokładnie wtedy, gdy
$p>1$. Rzeczywiście, podstawienie klasycznej
całki Riemanna $u=\log(1/x)$ na
$[e^{-R},e^{-1}]$ (a następnie twierdzenie
\ref{thm:mct}) daje
$\int_{[e^{-R},e^{-1}]}f_p\,\mathrm d\lambda_1(x)
=\int_1^R u^{-p}\,\mathrm d\lambda_1(u)$.
Całka po prawej ma skończoną granicę przy
$R\to\infty$ dokładnie wtedy, gdy $p>1$.
Samo użycie podstawienia w tym przykładzie
można też zastąpić oszacowaniami dyadycznymi;
pełne twierdzenie o zamianie zmiennych
udowodnimy później.
\end{przyklad}

\section{Całki zależne od parametru}

\begin{twierdzenie}[Ciągłość pod znakiem całki]
\label{thm:parameter-continuity}
Niech $T$ będzie przestrzenią metryczną, $t_0\in T$
i $F:T\times X\to\mathbb C$ funkcją taką, że
$x\mapsto F(t,x)$ jest mierzalna dla każdego $t\in T$.
Załóżmy, że istnieje otoczenie $V$ punktu $t_0$
oraz całkowalna $g:X\to[0,\infty]$ takie, że
$|F(t,x)|\le g(x)$ dla wszystkich $t\in V$ i
prawie wszystkich $x$, poza jednym wspólnym zbiorem
zerowym, oraz że $F(t,x)\to F(t_0,x)$, gdy
$t\to t_0$, dla prawie wszystkich $x$.
Wówczas funkcja
$\Phi(t)=\int_XF(t,x)\,\mathrm d\mu(x)$
jest określona dla $t\in V$ i ciągła w $t_0$.
\end{twierdzenie}
\begin{proof}
Ustalmy dowolny ciąg $t_n\to t_0$ w $T$.
Od pewnego indeksu $t_n\in V$. Funkcje
$F(t_n,\cdot)$ są zdominowane przez $g$ poza
wspólnym zbiorem zerowym i zbiegają prawie wszędzie
do $F(t_0,\cdot)$. Zespolona wersja twierdzenia
\ref{thm:dct} daje $\Phi(t_n)\to\Phi(t_0)$.
Ponieważ na przestrzeni metrycznej ciągowa
ciągłość jest równoważna ciągłości, otrzymujemy tezę.
\end{proof}

\begin{twierdzenie}[Różniczkowanie pod znakiem całki]
\label{thm:differentiating-integrals}
Niech $I\subseteq\mathbb R$ będzie przedziałem otwartym,
$t_0\in I$, a $F:I\times X\to\mathbb R$ funkcją,
dla której $x\mapsto F(t,x)$ jest mierzalna
dla każdego $t\in I$. Załóżmy, że:
\begin{enumerate}[label=\textup{(\roman*)}]
\item $F(t_0,\cdot)$ jest całkowalna;
\item istnieją $\delta>0$ o
$[t_0-\delta,t_0+\delta]\subset I$ i całkowalna
$g:X\to[0,\infty]$ takie, że dla prawie wszystkich
$x$ funkcja $t\mapsto F(t,x)$ jest różniczkowalna
na $(t_0-\delta,t_0+\delta)$ i
$|\partial_tF(t,x)|\le g(x)$ dla wszystkich $t$ z
tego przedziału (zbiór wyjątkowy jest wspólny).
\end{enumerate}
Wówczas $F(t,\cdot)$ jest całkowalna dla
$t$ dostatecznie bliskich $t_0$, funkcja
$\Phi(t)=\int_XF(t,x)\,\mathrm d\mu(x)$ jest
różniczkowalna w $t_0$ i
\begin{equation}\label{eq:differentiation-under-integral}
\Phi'(t_0)=\int_X\partial_tF(t_0,x)\,\mathrm d\mu(x).
\end{equation}
\end{twierdzenie}
\begin{proof}
Poza wspólnym zbiorem zerowym zastosujmy twierdzenie
o wartości średniej do funkcji $t\mapsto F(t,x)$.
Dla $|h|<\delta$ mamy
$|F(t_0+h,x)-F(t_0,x)|\le |h|g(x)$.
Zatem $|F(t_0+h,x)|\le|F(t_0,x)|+|h|g(x)$,
a prawa strona jest całkowalna. To dowodzi
istnienia całek w otoczeniu punktu $t_0$.
Funkcje mierzalne
$Q_h(x)=(F(t_0+h,x)-F(t_0,x))/h$ spełniają
$|Q_h(x)|\le g(x)$ dla wszystkich małych $h\ne0$
i zbiegają dla prawie wszystkich $x$ do
$\partial_tF(t_0,x)$. Granica punktowa ilorazów $Q_{1/n}$ jest mierzalna,
więc funkcja $x\mapsto\partial_tF(t_0,x)$,
po zdefiniowaniu jej jako zera na wspólnym zbiorze
wyjątkowym, jest mierzalna. Zbieżność zdominowana,
zastosowana do dowolnego ciągu $h_n\to0$,
$h_n\ne0$, daje
$\lim_n\int_XQ_{h_n}\,\mathrm d\mu(x)
=\int_X\partial_tF(t_0,x)\,\mathrm d\mu(x)$.
Z liniowości całki lewa strona jest granicą
ilorazów różnicowych funkcji $\Phi$; niezależność
od wybranego ciągu dowodzi istnienia pochodnej.
\end{proof}

\begin{przyklad}[Transformata Laplace'a funkcji ograniczonej]
Niech $u:[0,\infty)\to\mathbb R$ będzie mierzalna,
$|u(x)|\le M<\infty$, i dla $t>0$ określmy
$L(t)=\int_{[0,\infty)}e^{-tx}u(x)\,\mathrm d\lambda_1(x)$.
Całka jest skończona, gdyż $|e^{-tx}u(x)|\le Me^{-tx}$,
a całka z $e^{-tx}$ wynosi $1/t$ (z granicy
całek Riemanna na $[0,R]$, używając Beppo Leviego).
Dla ustalonego $t_0>0$ i $t\in(t_0/2,3t_0/2)$
mamy
$|\partial_t(e^{-tx}u(x))|=xe^{-tx}|u(x)|
\le Mxe^{-t_0x/2}$.
Funkcja dominująca jest całkowalna, np. z oszacowania
$x e^{-t_0x/2}\le C e^{-t_0x/4}$, gdzie $C<\infty$.
Twierdzenie \ref{thm:differentiating-integrals} daje
$L'(t)=-\int_{[0,\infty)}xe^{-tx}u(x)\,\mathrm d\lambda_1(x)$.
Analogicznie, przez indukcję, otrzymamy
$L^{(k)}(t)=(-1)^k\int_{[0,\infty)}x^ke^{-tx}u(x)
\,\mathrm d\lambda_1(x)$ dla wszystkich $k\ge1$.
Istotnie, dla ustalonych $t_0>0$ i $k\ge0$
wszystkie potrzebne pochodne są lokalnie
zdominowane przez $Mx^{k+1}e^{-t_0x/2}$;
funkcja ta jest całkowalna, ponieważ dla
$d=t_0/4$ iloczyn $x^{k+1}e^{-dx}$ jest
ograniczony na $[0,\infty)$.
\end{przyklad}

\section{Jednostajna całkowalność i zbieżność według miary}

\begin{definicja}[Jednostajna całkowalność]
\label{def:uniform-integrability}
Niech $\mu(X)<\infty$. Niepustą rodzinę $\mathcal U$
mierzalnych, całkowalnych funkcji rzeczywistych
na $X$ nazywamy \emph{jednostajnie całkowalną},
jeśli
\[
\lim_{K\to\infty}\sup_{u\in\mathcal U}
\int_{\{|u|>K\}}|u(x)|\,\mathrm d\mu(x)=0.
\]
Warunek kontroluje masę części funkcji o dużych
wartościach, a nie jedynie wielkość miary poziomic.
\end{definicja}

\begin{lemat}[Kontrola całek na małych zbiorach]
\label{lem:uniform-absolute-continuity}
Jeśli $\mu(X)<\infty$ i $\mathcal U$ jest jednostajnie
całkowalną rodziną funkcji rzeczywistych, to:
\begin{enumerate}[label=\textup{(\roman*)}]
\item $\sup_{u\in\mathcal U}\int_X|u|\,\mathrm d\mu(x)<\infty$;
\item dla każdego $\varepsilon>0$ istnieje $\delta>0$
takie, że dla każdego $A\in\mathcal F$ z
$\mu(A)<\delta$ zachodzi
$\sup_{u\in\mathcal U}\int_A|u|\,\mathrm d\mu(x)<\varepsilon$.
\end{enumerate}
\end{lemat}
\begin{proof}
Ustalmy $\varepsilon>0$ i wybierzmy $K>0$ takie, że
$\sup_{u\in\mathcal U}\int_{\{|u|>K\}}|u|<\varepsilon/2$.
Dla każdego $u\in\mathcal U$ i $A\in\mathcal F$
z podziału na $\{|u|\le K\}$ i $\{|u|>K\}$
wynika
$\int_A|u|\,\mathrm d\mu(x)\le K\mu(A)+\varepsilon/2$.
Biorąc $A=X$, dowodzimy (i) przy skończonej
mierze $\mu(X)$. Wybór
$\delta=\varepsilon/(2K)$ dla $K>0$ daje (ii);
przy $K=0$ oszacowanie jest nawet prostsze.
\end{proof}

\begin{twierdzenie}[Dwa równoważne opisy jednostajnej całkowalności]
\label{thm:ui-two-criteria-8}
Niech $\mu(X)<\infty$ i niech $\mathcal U$ będzie
rodziną rzeczywistych funkcji mierzalnych $u:X\to\mathbb R$
spełniających $\int_X|u|\,\mathrm d\mu(x)<\infty$.
Równość funkcji prawie wszędzie nie ma wpływu
na żadną z rozważanych własności.
Wówczas $\mathcal U$ jest jednostajnie całkowalna
w sensie definicji~\ref{def:uniform-integrability}
wtedy i tylko wtedy, gdy jednocześnie:
\begin{enumerate}[label=\textup{(\roman*)}]
\item $M:=\sup_{u\in\mathcal U}\int_X|u|\,\mathrm d\mu(x)<\infty$;
\item dla każdego $\varepsilon>0$ istnieje $\delta>0$
takie, że $\mu(A)<\delta$ implikuje
$\sup_{u\in\mathcal U}\int_A|u|\,\mathrm d\mu(x)<\varepsilon$.
\end{enumerate}
\end{twierdzenie}
\begin{proof}
Konieczność to lemat~\ref{lem:uniform-absolute-continuity}.
Odwrotnie, przy $M=0$ wszystkie funkcje są zerowe
prawie wszędzie, więc teza jest natychmiastowa.
Dla $M>0$ i zadanego $\varepsilon>0$ wybieramy
$\delta$ z (ii). Nierówność Markowa daje
$\mu(\{|u|>K\})\le M/K$ dla każdego $u\in\mathcal U$.
Gdy $K>M/\delta$, zbiór $\{|u|>K\}$ ma miarę
mniejszą niż $\delta$, zatem (ii) daje
$\int_{\{|u|>K\}}|u|<\varepsilon$ jednostajnie
względem $u$. Jest to dokładnie warunek z
definicji~\ref{def:uniform-integrability}.
\end{proof}

\begin{twierdzenie}[Zbieżność Vitaliego na przestrzeni o mierze skończonej]
\label{thm:vitali-l1}
Niech $\mu(X)<\infty$, a $f_n:X\to\mathbb R$ będą
mierzalnymi funkcjami tworzącymi rodzinę
jednostajnie całkowalną. Jeśli
$f_n\xrightarrow{\mu}f$ dla pewnej mierzalnej
$f:X\to\mathbb R$, to $f$ jest całkowalna oraz
\[
\lim_{n\to\infty}\int_X|f_n(x)-f(x)|\,\mathrm d\mu(x)=0.
\]
W szczególności $\int_Xf_n\,\mathrm d\mu(x)\to
\int_Xf\,\mathrm d\mu(x)$.
\end{twierdzenie}
\begin{proof}
Z twierdzenia \ref{thm:measure-subsequence}
można wybrać podciąg $f_{n_k}\to f$ prawie wszędzie.
Na mocy lematu Fatou i części (i) lematu
\ref{lem:uniform-absolute-continuity}
$\int_X|f|\,\mathrm d\mu(x)\le
\liminf_k\int_X|f_{n_k}|\,\mathrm d\mu(x)<\infty$.

Ustalmy $\varepsilon>0$. Z lematu
\ref{lem:uniform-absolute-continuity} wybieramy
$\delta_1>0$ takie, że
$\mu(A)<\delta_1\Rightarrow
\sup_n\int_A|f_n|\,\mathrm d\mu(x)<\varepsilon/3$.
Ponieważ $f$ jest całkowalna, również jej całka
jest absolutnie ciągła względem miary w następującym
sensie: istnieje $\delta_2>0$ takie, że
$\mu(A)<\delta_2\Rightarrow\int_A|f|<\varepsilon/3$.
Istotnie, wystarczy obciąć $|f|$ na poziomie $K$,
tak by $\int_{\{|f|>K\}}|f|<\varepsilon/6$,
i oszacować pozostałą część przez $K\mu(A)$.

Niech $\eta>0$ będzie tak małe, aby
$\eta\mu(X)<\varepsilon/3$. Ze zbieżności według
miary wynika, że dla dużych $n$ zbiór
$A_n=\{|f_n-f|>\eta\}$ ma miarę mniejszą od
$\min(\delta_1,\delta_2)$. Dzieląc całkę na
$A_n$ i dopełnienie, dostajemy
\[
\int_X|f_n-f|\,\mathrm d\mu(x)
\le\eta\mu(X)+\int_{A_n}|f_n|\,\mathrm d\mu(x)
+\int_{A_n}|f|\,\mathrm d\mu(x)<\varepsilon.
\]
To dowodzi zbieżności w całce z modułu różnicy.
Równość granic całek wynika z oszacowania
$|\int_X f_n-\int_X f|\le\int_X|f_n-f|$.
\end{proof}

\begin{przyklad}[Masę można skoncentrować na małych zbiorach]
Na $(0,1)$ niech $f_n=n\chi_{(0,1/n)}$.
Ciąg zbiega do zera według miary, ponieważ
$\lambda_1(\{|f_n|>\varepsilon\})\le1/n$
dla dużych $n$. Rodzina nie jest jednak jednostajnie
całkowalna: dla każdego $K>0$ i każdego $n>K$
zachodzi
$\int_{\{|f_n|>K\}}|f_n|\,\mathrm d\lambda_1(x)=1$.
Twierdzenie Vitaliego zatem nie ma zastosowania.
\end{przyklad}

\section{Związek z całką Riemanna}

\begin{twierdzenie}[Zgodność całki Riemanna i Lebesgue'a]
\label{thm:riemann-to-lebesgue}
Niech $a<b$ i $f:[a,b]\to\mathbb R$ będzie funkcją
ograniczoną, całkowalną w sensie Riemanna.
Wówczas $f$ jest mierzalna względem uzupełnionej
miary Lebesgue'a na $[a,b]$, jest całkowalna
w sensie Lebesgue'a, a obie całki są równe:
\[
\int_a^b f(x)\,\mathrm dx
=\int_{[a,b]}f(x)\,\mathrm d\lambda_1(x).
\]
Po lewej stronie występuje tradycyjna całka Riemanna,
a po prawej całka względem miary Lebesgue'a.
\end{twierdzenie}
\begin{proof}
Niech $M$ będzie ograniczeniem $|f|$.
Z kryterium Darboux dla całkowalności Riemanna
możemy wybrać podziały $P_n$ takie, że
$U(f,P_n)-L(f,P_n)<2^{-2n}$,
gdzie $U,L$ są odpowiednio górnymi i dolnymi
sumami Darboux. Z każdego podziału budujemy
proste funkcje borelowskie $u_n,\ell_n$:
na rozłącznych półotwartych częściach podziału
przypisujemy odpowiednio supremum i infimum $f$
na odpowiadających im \emph{domkniętych} przedziałach
podziału (ostatni przedział zawiera prawy koniec).
Dzięki temu $\ell_n\le f\le u_n$ w każdym punkcie,
a $\int_{[a,b]}u_n(x)\,\mathrm d\lambda_1(x)=U(f,P_n)$,
$\int_{[a,b]}\ell_n(x)\,\mathrm d\lambda_1(x)=L(f,P_n)$.

Ponieważ $u_n-\ell_n\ge0$ i
$\int_{[a,b]}(u_n-\ell_n)\,\mathrm d\lambda_1(x)
<2^{-2n}$, nierówność Markowa daje
$\lambda_1(\{u_n-\ell_n\ge2^{-n}\})\le2^{-n}$.
Pierwsza część lematu Borela--Cantellego implikuje,
że dla prawie wszystkich $x$ istnieje $N(x)$ takie,
że $u_n(x)-\ell_n(x)<2^{-n}$ dla $n\ge N(x)$.
Z nierówności $\ell_n\le f\le u_n$ otrzymujemy
$u_n(x)\to f(x)$ poza mierzalnym zbiorem zerowym.
Granica mierzalnych funkcji jest mierzalna poza
tym zbiorem; dzięki zupełności miary Lebesgue'a
sama $f$ jest mierzalna. Ponadto $|f|\le M$
i $\lambda_1([a,b])<\infty$, więc $f$ jest całkowalna.

Z monotoniczności całki i równości dla prostych
funkcji otrzymujemy
$L(f,P_n)\le\int_{[a,b]}f\,\mathrm d\lambda_1(x)
\le U(f,P_n)$. Obie skrajne liczby dążą do
całki Riemanna, bo ich różnica dąży do zera,
a całka Riemanna leży pomiędzy nimi.
Twierdzenie o trzech ciągach kończy dowód.
\end{proof}

\begin{twierdzenie}[Kryterium Lebesgue'a całkowalności Riemanna]
\label{thm:riemann-criterion}
Niech $f:[a,b]\to\mathbb R$ będzie ograniczona, gdzie
$a<b$. Wówczas $f$ jest całkowalna w sensie Riemanna
tedy i tylko wtedy, gdy zbiór punktów nieciągłości $f$
ma miarę Lebesgue'a zero.
\end{twierdzenie}
\begin{proof}
Dla $x\in[a,b]$ oznaczmy przez $\omega_f(x)$
oscylację $f$ w punkcie, czyli
\[
\omega_f(x):=\inf_{r>0}
\left(\sup_{y\in[a,b]\cap(x-r,x+r)}f(y)
-\inf_{y\in[a,b]\cap(x-r,x+r)}f(y)\right).
\]
Funkcja $f$ jest ciągła w $x$ wtedy i tylko wtedy,
gdy $\omega_f(x)=0$. Dla $\varepsilon>0$ zbiór
$D_\varepsilon=\{x:\omega_f(x)\ge\varepsilon\}$
jest domknięty: jeśli $\omega_f(x)<\varepsilon$,
to na pewnym otoczeniu punktu $x$ oscylacja $f$
jest mniejsza niż $\varepsilon$, a zatem także
oscylacje w punktach pobliskich są mniejsze od
$\varepsilon$. Zatem $D_\varepsilon$ jest zwarty.

Załóżmy najpierw całkowalność Riemanna i
ustalmy $\varepsilon>0$. Weźmy dowolny podział
$P=(a=x_0<\cdots<x_N=b)$. Jeśli
$x\in D_\varepsilon$ nie jest punktem podziału,
to oscylacja $f$ na domkniętym przedziale
$[x_{j-1},x_j]$ zawierającym $x$ jest co najmniej
$\varepsilon$. Dlatego $D_\varepsilon$ mieści się
w sumie skończenie wielu punktów podziału i tych
przedziałów, których oscylacja jest nie mniejsza
niż $\varepsilon$. Suma długości tych przedziałów
jest nie większa niż
$\varepsilon^{-1}(U(f,P)-L(f,P))$.
Ponieważ dla funkcji całkowalnej w sensie Riemanna
różnicę tę można uczynić dowolnie małą,
$\lambda_1(D_\varepsilon)=0$.
Zbiór nieciągłości jest sumą
$\bigcup_{m\ge1}D_{1/m}$, więc jest zerowy.

Odwrotnie, niech zbiór nieciągłości $D$ ma miarę
zero, a $|f|\le M$. Ustalmy $\varepsilon>0$.
Zbiór $D_\varepsilon$ jest zwarty i zerowy, więc
istnieje zbiór otwarty $O$ zawierający $D_\varepsilon$
o dowolnie małej mierze Lebesgue'a (regularność
zewnętrzna, rozdział 4). Zbiór
$K=[a,b]\setminus O$ jest zwarty. Dla każdego
$x\in K$ mamy $\omega_f(x)<\varepsilon$,
zatem istnieje promień $r_x>0$, na którego kuli
$B(x,r_x)\cap[a,b]$ oscylacja $f$ jest
mniejsza niż $\varepsilon$. Zwartość $K$ pozwala
wybrać skończenie wiele kul $B(x_j,r_{x_j}/2)$
pokrywających $K$. Połóżmy
$\delta=\frac12\min_j r_{x_j}>0$.
Jeżeli domknięty przedział $J\subseteq[a,b]$
o długości mniejszej niż $\delta$ przecina $K$,
wybieramy $y\in J\cap K$ i indeks $j$ taki, że
$y\in B(x_j,r_{x_j}/2)$. Dla każdego $z\in J$
mamy
$|z-x_j|\le|z-y|+|y-x_j|<\delta+r_{x_j}/2
\le r_{x_j}$.
Zatem $J\subseteq B(x_j,r_{x_j})$, gdzie
oscylacja $f$ jest mniejsza niż $\varepsilon$.

Weźmy podział $P$ o siatce mniejszej niż $\delta$.
Na przedziałach przecinających $K$ oscylacja $f$
jest mniejsza niż $\varepsilon$. Pozostałe
przedziały (z wyjątkiem ewentualnych końców)
mieszczą się w $O$, więc suma ich długości
nie przekracza $\lambda_1(O)$. Otrzymujemy
\[
U(f,P)-L(f,P)
\le\varepsilon(b-a)+2M\lambda_1(O).
\]
Wybierając najpierw $\varepsilon$, a potem
$\lambda_1(O)$ dostatecznie małe, czynimy prawą
stronę dowolnie małą. Kryterium Darboux dowodzi
całkowalności Riemanna.
\end{proof}

\begin{uwaga}[Dlaczego oba warunki mają znaczenie?]
Funkcja Dirichleta $\chi_{\mathbb Q\cap[0,1]}$
jest całkowalna w sensie Lebesgue'a, ale
nie jest całkowalna w sensie Riemanna, gdyż
zbiór jej punktów nieciągłości to cały $[0,1]$.
Z kolei każda ograniczona funkcja o przeliczalnym
zbiorze punktów nieciągłości jest całkowalna
w sensie Riemanna. Warunek ograniczoności jest
w kryterium konieczny w jego klasycznej postaci:
na przykład $x\mapsto1/x$ na $(0,1]$ nie jest
ograniczona, mimo ciągłości w całej swojej dziedzinie.
\end{uwaga}

\section{Podsumowanie logiczne}

Twierdzenie Beppo Leviego opiera się na definicji
całki jako supremum po funkcjach prostych i na
ciągłości miary z dołu. Z niego wynikają
addytywność całki nieujemnej i możliwość całkowania
szeregów nieujemnych. Lemat Fatou to zastosowanie
Beppo Leviego do ciągu infimów ogonowych, a twierdzenie
o zbieżności zdominowanej uzyskujemy z lematu Fatou
przez rozważenie nieujemnych funkcji $g\pm f_n$.
Wszystkie te rezultaty pozostają ważne dla miar
nieskończonych; wyjątkiem jest twierdzenie o zbieżności
ograniczonej, w którym skończoność miary
zapewnia całkowalność stałej funkcji dominującej.